% Options for packages loaded elsewhere
\PassOptionsToPackage{unicode}{hyperref}
\PassOptionsToPackage{hyphens}{url}
\documentclass[
]{article}
\usepackage{xcolor}
\usepackage[margin=1in]{geometry}
\usepackage{amsmath,amssymb}
\setcounter{secnumdepth}{-\maxdimen} % remove section numbering
\usepackage{iftex}
\ifPDFTeX
  \usepackage[T1]{fontenc}
  \usepackage[utf8]{inputenc}
  \usepackage{textcomp} % provide euro and other symbols
\else % if luatex or xetex
  \usepackage{unicode-math} % this also loads fontspec
  \defaultfontfeatures{Scale=MatchLowercase}
  \defaultfontfeatures[\rmfamily]{Ligatures=TeX,Scale=1}
\fi
\usepackage{lmodern}
\ifPDFTeX\else
  % xetex/luatex font selection
\fi
% Use upquote if available, for straight quotes in verbatim environments
\IfFileExists{upquote.sty}{\usepackage{upquote}}{}
\IfFileExists{microtype.sty}{% use microtype if available
  \usepackage[]{microtype}
  \UseMicrotypeSet[protrusion]{basicmath} % disable protrusion for tt fonts
}{}
\makeatletter
\@ifundefined{KOMAClassName}{% if non-KOMA class
  \IfFileExists{parskip.sty}{%
    \usepackage{parskip}
  }{% else
    \setlength{\parindent}{0pt}
    \setlength{\parskip}{6pt plus 2pt minus 1pt}}
}{% if KOMA class
  \KOMAoptions{parskip=half}}
\makeatother
\usepackage{longtable,booktabs,array}
\usepackage{calc} % for calculating minipage widths
% Correct order of tables after \paragraph or \subparagraph
\usepackage{etoolbox}
\makeatletter
\patchcmd\longtable{\par}{\if@noskipsec\mbox{}\fi\par}{}{}
\makeatother
% Allow footnotes in longtable head/foot
\IfFileExists{footnotehyper.sty}{\usepackage{footnotehyper}}{\usepackage{footnote}}
\makesavenoteenv{longtable}
\usepackage{graphicx}
\makeatletter
\newsavebox\pandoc@box
\newcommand*\pandocbounded[1]{% scales image to fit in text height/width
  \sbox\pandoc@box{#1}%
  \Gscale@div\@tempa{\textheight}{\dimexpr\ht\pandoc@box+\dp\pandoc@box\relax}%
  \Gscale@div\@tempb{\linewidth}{\wd\pandoc@box}%
  \ifdim\@tempb\p@<\@tempa\p@\let\@tempa\@tempb\fi% select the smaller of both
  \ifdim\@tempa\p@<\p@\scalebox{\@tempa}{\usebox\pandoc@box}%
  \else\usebox{\pandoc@box}%
  \fi%
}
% Set default figure placement to htbp
\def\fps@figure{htbp}
\makeatother
\setlength{\emergencystretch}{3em} % prevent overfull lines
\providecommand{\tightlist}{%
  \setlength{\itemsep}{0pt}\setlength{\parskip}{0pt}}
\usepackage{bookmark}
\IfFileExists{xurl.sty}{\usepackage{xurl}}{} % add URL line breaks if available
\urlstyle{same}
\hypersetup{
  hidelinks,
  pdfcreator={LaTeX via pandoc}}

\author{}
\date{\vspace{-2.5em}}

\begin{document}

\section{Data Analysis Plan and Power Justification: Aims 1 and
2}\label{data-analysis-plan-and-power-justification-aims-1-and-2}

\subsection{C.5.e. Aim 1 Analytical
Plan}\label{c.5.e.-aim-1-analytical-plan}

\textbf{Data Handling.} Distributional properties of cytokine/chemokine
variables will be examined via Q-Q plots and Shapiro-Wilk tests;
log-transformation will be applied only if warranted. Should any values
fall below the assay detection limit, they will be substituted at half
that limit. Outliers exceeding 4 SD will be flagged. Primary analyses
will use available cases under a missing-at-random assumption; multiple
imputation and inverse-probability-of-attrition weighting will serve as
sensitivity analyses if dropout is associated with baseline
characteristics.

\textbf{Analytic Approach.} The primary analysis will use multivariable
linear regression with the one-year change score (follow-up minus
baseline) as the dependent variable and the corresponding baseline
outcome value as a covariate to adjust for regression to the mean. For
H1a, the primary predictor is the baseline cytokine/chemokine level; for
H1b, the change in that marker. Separate models will be fit for each of
four primary markers (IL-6, TNF-α, MCP-1, Eotaxin-1) across two
outcomes: (1) change in episodic memory consolidation (CVLT, pattern
separation, long-term consolidation composite) and (2) change in
AD-signature mean cortical thickness (FreeSurfer 5.1). Markers will not
be combined as simultaneous predictors due to high inter-marker
correlation (preliminary IL-6/MCP-1: \emph{r} = 0.93). All models will
adjust for age, sex, APOE ε4 status, BMI, hypercholesterolemia history,
NSAID use, immune-related health conditions, and diagnostic group (aMCI
vs.~HC). Both groups are combined to maximize the range of cytokine and
outcome levels; effect modification by group will be examined
secondarily. Collinearity will be assessed via VIF (\textgreater5
triggers investigation); covariates will also be evaluated using the
change-in-estimate criterion (\textgreater10\% change in primary
estimate warrants retention). Total cortical thickness will be included
in MRI outcome models to isolate AD-signature effects. Table A1
(Appendix) summarizes all variables by role.

\subsection{C.6.c. Aim 2 Analytical
Plan}\label{c.6.c.-aim-2-analytical-plan}

Aim 2 extends the Aim 1 framework by adding baseline Florbetapir-PET
amyloid deposition (global SUVR) and its product interaction with each
cytokine/chemokine. The primary test of H2 is the interaction
coefficient, evaluating whether co-occurring elevated inflammation and
amyloid burden predicts steeper decline than either factor alone. The
same covariates and outcomes from Aim 1 apply; both main effects are
included. As a secondary analysis, amyloid will be dichotomized at SUVR
≥ 1.10; a significant interaction will be probed via simple slopes (±1
SD of cytokine), yielding a 2×2 interpretation:
amyloid-positive/high-inflammation, amyloid-positive/low-inflammation,
amyloid-negative/high-inflammation, and
amyloid-negative/low-inflammation. The hypothesis predicts the steepest
decline in the amyloid-positive/high-inflammation cell. Sensitivity
analyses stratified by APOE ε4 will evaluate whether the interaction
differs by genotype.

\textbf{Multiple Comparisons.} Within each outcome domain, false
discovery rate (FDR; Benjamini-Hochberg, q \textless{} 0.05) will be
applied across the four inflammatory markers. FDR is preferred over
Bonferroni because the markers are substantially correlated, rendering
family-wise error control overly conservative. Bonferroni correction
(α/4 = 0.0125) will be evaluated as a sensitivity analysis. The two
outcome domains (memory, cortical thickness; preliminary \emph{r} =
0.16) will not be corrected across. Concordant patterns across
correlated markers (e.g., both MCP-1 and Eotaxin-1 predicting decline)
will strengthen evidence for an inflammatory pathway even if individual
tests do not survive strict correction.

\subsection{D. Power Analysis and Sample Size
Justification}\label{d.-power-analysis-and-sample-size-justification}

We estimated power for N = 175 (125 aMCI + 50 HC, after 10\% attrition
from 192 enrolled). Effect sizes were derived from Pearson correlations
in the investigator's preliminary dataset (n = 30), ranging from
\textbar{}\emph{r}\textbar{} = 0.26 (IL-6 vs.~memory change) to
\textbar{}\emph{r}\textbar{} = 0.69 (MCP-1 vs.~cortical thickness
change). All tests were two-sided.

\textbf{Aim 1.} Power was estimated using \texttt{pwr.r.test()} (R
\texttt{pwr} package), approximating the regression test as a bivariate
correlation --- a conservative approach since covariate adjustment
generally increases precision. At N = 175, the minimum detectable
\textbar{}\emph{r}\textbar{} at 80\% power was 0.21 (α = 0.05) and 0.25
(α = 0.0125). The weakest preliminary signal (IL-6/memory,
\textbar{}\emph{r}\textbar{} = 0.26) achieved 94\% and 84\% power at
these thresholds, respectively; all stronger associations exceeded 96\%
(Table 1, Figure 1A). Because the effective FDR-adjusted alpha falls
between these bounds, the study is well-powered for Aim 1.

\textbf{Aim 2.} Power for the interaction was estimated using
\texttt{pwr.f2.test()} (R \texttt{pwr} package), testing the R²
increment of the interaction beyond main effects and covariates (10
total predictors; numerator df = 1, denominator df = 164). The minimum
detectable interaction at 80\% power was f² = 0.048 (α = 0.05) and f² =
0.069 (α = 0.0125), both small-to-medium by Cohen's benchmarks. At f² =
0.05, power was 82\% (α = 0.05); at f² = 0.08, power reached 87\% even
at α = 0.0125 (Table 2, Figure 1B). Assuming the full model explains
\textasciitilde30\% of outcome variance, the detectable interaction (f²
= 0.048) corresponds to ΔR² ≈ 3.4\% --- a clinically meaningful
acceleration of decline when elevated inflammation co-occurs with
amyloid positivity. If the investigator's enrollment capacity permits,
increasing N would further improve power for the interaction at
corrected alpha levels.

\begin{center}\rule{0.5\linewidth}{0.5pt}\end{center}

\textbf{Table 1. Aim 1 Power: Bivariate Correlation Test (N = 175,
two-sided)}

\begin{longtable}[]{@{}lll@{}}
\toprule\noalign{}
& \emph{r} & \\
\midrule\noalign{}
\endhead
\bottomrule\noalign{}
\endlastfoot
0.10 & 0.261 & 0.119 \\
0.15 & 0.511 & 0.304 \\
0.20 & 0.760 & 0.565 \\
0.25 & 0.919 & 0.804 \\
0.30 & 0.983 & 0.942 \\
0.35 & 0.998 & 0.989 \\
0.40 & 1.000 & 0.999 \\
0.45 & 1.000 & 1.000 \\
0.50 & 1.000 & 1.000 \\
\end{longtable}

\emph{\texttt{pwr.r.test()}, R \texttt{pwr} package. Preliminary
\textbar r\textbar{} ranged 0.26--0.69.}

\begin{center}\rule{0.5\linewidth}{0.5pt}\end{center}

\textbf{Table 2. Aim 2 Power: Interaction F-test (N = 175, u = 1, v =
164)}

\begin{longtable}[]{@{}llll@{}}
\toprule\noalign{}
f² & ΔR²* & Power (α = 0.05) & Power (α = 0.0125) \\
\midrule\noalign{}
\endhead
\bottomrule\noalign{}
\endlastfoot
0.02 & 0.014 & 0.441 & 0.244 \\
0.03 & 0.021 & 0.602 & 0.387 \\
0.04 & 0.028 & 0.726 & 0.522 \\
0.05 & 0.035 & 0.817 & 0.639 \\
0.06 & 0.042 & 0.880 & 0.735 \\
0.08 & 0.056 & 0.952 & 0.867 \\
0.10 & 0.070 & 0.982 & 0.938 \\
0.15 & 0.105 & 0.999 & 0.993 \\
\end{longtable}

\emph{\texttt{pwr.f2.test()}, R \texttt{pwr} package. *ΔR² ≈ f² × (1 −
R²\_full), assuming R²\_full ≈ 0.30. Cohen: small = 0.02, medium =
0.15.}

\begin{center}\rule{0.5\linewidth}{0.5pt}\end{center}

\textbf{Table A1. Variable Summary by Role}

\begin{longtable}[]{@{}
  >{\raggedright\arraybackslash}p{(\linewidth - 6\tabcolsep) * \real{0.1364}}
  >{\raggedright\arraybackslash}p{(\linewidth - 6\tabcolsep) * \real{0.2273}}
  >{\raggedright\arraybackslash}p{(\linewidth - 6\tabcolsep) * \real{0.1364}}
  >{\raggedright\arraybackslash}p{(\linewidth - 6\tabcolsep) * \real{0.5000}}@{}}
\toprule\noalign{}
\begin{minipage}[b]{\linewidth}\raggedright
Role
\end{minipage} & \begin{minipage}[b]{\linewidth}\raggedright
Variable
\end{minipage} & \begin{minipage}[b]{\linewidth}\raggedright
Type
\end{minipage} & \begin{minipage}[b]{\linewidth}\raggedright
Measurement / Coding
\end{minipage} \\
\midrule\noalign{}
\endhead
\bottomrule\noalign{}
\endlastfoot
Primary Outcome (Both Aims) & Change in episodic memory consolidation &
Continuous & CVLT + pattern separation + LT consolidation; follow-up −
baseline \\
Primary Outcome (Both Aims) & Change in AD-signature cortical thickness
& Continuous & FreeSurfer 5.1 AD-signature ROI; follow-up − baseline \\
Primary Predictor (Aim 1) & Plasma cytokines/chemokines (IL-6, TNF-α,
MCP-1, Eotaxin-1) & Continuous & pg/mL; transform if warranted; H1a:
baseline; H1b: change \\
Primary Predictor (Aim 2) & Amyloid SUVR × cytokine interaction &
Continuous; Binary & Primary: SUVR × cytokine product; Secondary:
amyloid-pos/neg (≥1.10) × cytokine \\
Covariates (All Models) & Age, sex, APOE ε4, BMI, hypercholesterolemia,
NSAID use, immune conditions, diagnostic group, baseline outcome &
Continuous; Binary & Pre-specified; baseline outcome adjusts for
regression to the mean \\
\end{longtable}

\end{document}
